The fresh LLM mean is 8.23\%, versus 9.36\% for the rule and 9.69\% for random selection. The paired LLM-minus-rule difference is -1.12 percentage points (pp), and LLM-minus-random is -1.46 pp. Exact case-level sign-flip $p$ values are 0.379 and 0.405; city-level values are 0.188 and 0.688. Holm-adjusted Wilcoxon values are 0.346 and 0.347. The city-bootstrap intervals are [-2.53, -0.01] pp and [-5.08, 0.95] pp. The first descriptive bootstrap interval narrowly excludes zero, whereas the discrete city sign-flip test does not reject; with five cities these should not be read as precise confirmatory bounds. No comparison supports an LLM advantage. The LLM wins/ties/loses 2/6/7 cases against the rule and 5/3/7 against random.

The LLM produces 10 failed refinement decisions out of 30, compared with none for the non-LLM selectors. Its 92 candidate SUMO invocations yield 91 successful simulations; the other arms each use 103 invocations with 102 successful simulations. Pooled feasible-improving yields are 31.5\%, 38.8\%, and 36.9\%, respectively. Thus a shared budget cap does not imply equal realized traffic evaluations. We report this reliability cost, rather than replacing missing LLM choices with a stronger policy.
